\begin{initiativkarte}{RobOtter Club Hamburg (RCHH)}

  % Bild (optional)
  \IfFileExists{bilder/robotter_club.jpg}{%
    \begin{center}
      \includegraphics[width=\linewidth,height=5cm,keepaspectratio]{bilder/robotter_club.jpg}
    \end{center}
    \vspace{0.4cm}
  }{}

  % Kurzbeschreibung
  \textit{\textcolor{hawgray}{Der RobOtter Club Hamburg entwickelt autonome Roboter und nimmt regelmäßig an internationalen Robotik-Wettbewerben wie Eurobot teil.}}

  \vspace{0.5cm}
  \hawrule

  % Beschreibung
  \textbf{\textcolor{hawblue}{Was ist der RobOtter Club?}}\\[0.3cm]

  Der \textbf{RobOtter Club Hamburg} ist eine seit 2011 aktive studentische Initiative der HAW Hamburg im Bereich Robotik und autonome Systeme.

  Ziel des Clubs ist die regelmäßige Teilnahme an internationalen Wettbewerben wie \textbf{Eurobot}, bei denen vollständig selbst entwickelte Roboter gegeneinander antreten.

  \vspace{0.3cm}

  Dabei werden alle Komponenten selbst entwickelt – von Mechanik über Elektronik bis zur Software. Studierende wenden hier ihr im Studium erlerntes Wissen praxisnah an und vertiefen es in interdisziplinären Teams.

  \vspace{0.5cm}

  \textbf{Worum geht es im Club?}

  \begin{itemize}
    \item Entwicklung autonomer Roboter
    \item Teilnahme an internationalen Wettbewerben (Eurobot u.a.)
    \item Anwendung von Robotik, KI und Mechatronik
    \item Projektmanagement und Teamarbeit
    \item Zusammenarbeit mit Industrie und Forschung
  \end{itemize}

  \vspace{0.5cm}

  \textbf{Erfolge \& Geschichte}

  \begin{itemize}
    \item 2011: Erste Teilnahme an der Eurobot (Otter 1)
    \item 2013: Eurobot-Vorentscheid an der HAW Hamburg
    \item 2017: Deutscher Vizemeister und internationale Finalteilnahme
    \item 2018: Platz 18 beim internationalen Finale in Frankreich
  \end{itemize}

  \vspace{0.5cm}

  Neben dem Wettbewerb vermittelt der Club auch wichtige Kompetenzen in den Bereichen:

  \begin{itemize}
    \item Embedded Systems
    \item Konstruktion und Fertigung
    \item Softwareentwicklung
    \item Marketing und Organisation
  \end{itemize}

  \vspace{0.5cm}

  % Tags
  \tag{Robotik}
  \tag{Engineering}
  \tag{Wettbewerb}
  \tag{Mechatronik}
  \tag{Innovation}

  \vspace{0.6cm}
  \hawrule

  % Infozeilen
  \infozeile{\faMapMarker*}{Ort: Berliner Tor 21, Raum 224}
  \infozeile{\faUser}{Mentor: Prof. Dr. Thomas Frischgesell}

\end{initiativkarte}